\originalchapter{星系观测如何变成可比较的物理量}

本章对应 8.902 Fall 2023 第1讲与第2讲开头的恒星种群说明. 课程由 Mark Vogelsberger 主讲, 笔记由 Stephanie O'Neil 与 David DePalma 整理. 我们从探测器实际收到的光出发, 依次建立通量、星等、面亮度、光度函数和速度的联系. 多波段照片中的颜色只是观测量的一种显示, 要讨论恒星年龄或暗物质, 还须说明通带、距离与动力学模型. 本章不复制原讲义中第三方照片; 图均由给定数学模型自绘.

\originalsection{通量密度与比强度的区别}
设一束频率在 $[\nu,\nu+\dd\nu]$ 内的辐射穿过面积 $\dd A$. 面的法线与光传播方向夹角为 $\theta$, 比强度 $I_\nu$ 按
\[
\dd E=I_\nu\cos\theta\,\dd A\,\dd t\,\dd\nu\,\dd\Omega
\]
定义. 它把方向也分辨出来; 通量密度则对方向积分:
\[
F_\nu=\int I_\nu\cos\theta\dd\Omega.
\]
单位分别为 $\mathrm{erg\,s^{-1}cm^{-2}Hz^{-1}sr^{-1}}$ 与 $\mathrm{erg\,s^{-1}cm^{-2}Hz^{-1}}$. 对一个小角度的远方天体, 探测面朝向它, 可近似 $F_\nu=\int_{\rm source}I_\nu\dd\Omega$. 球对称、无吸收且静态的源若光谱光度为 $L_\nu$, 穿过任意包围它的球面的功率相同, 因此
\[
4\pi D^2F_\nu=L_\nu,\qquad F_\nu\propto D^{-2}.
\]
原笔记先列出的 $r^2$ 与随后能量守恒推导不一致; 此处以守恒式确定负二次方. 同一小块物理发光面的张角也随 $D^{-2}$ 变化, 故静态 Euclidean 条件下其比强度不因离远而改变. 总光变暗与每单位张角的光不变并不矛盾: 远处的源在图上占据更少的面积.

这条面亮度性质需附条件. 吸收、分辨率不足、宇宙红移都能改变测量. 后文将得到光子传播中的 $I_\nu/\nu^3$ 不变量, 由它推出宇宙学总面亮度的 $(1+z)^{-4}$ 衰减; 不能把这里的静态结论无条件用于高红移星系.

\originalsection{星等以对数记录通量比}
在同一通带和零点定义下,
\[
m_1-m_2=-2.5\log_{10}\frac{F_1}{F_2}.
\]
所以低星等表示更亮, 五个星等对应一百倍通量. Vega 系统以织女星的参考光谱定零点; AB 系统使用频率通量密度恒定的参考源 $F_{\nu,0}=3631\,\mathrm{Jy}$, $1\,\mathrm{Jy}=10^{-23}\,\mathrm{erg\,s^{-1}cm^{-2}Hz^{-1}}$. 单色 AB 星等写成
\[
m_{\rm AB}=-2.5\log_{10}\!\left[\frac{F_\nu}{\mathrm{erg\,s^{-1}cm^{-2}Hz^{-1}}}\right]-48.60.
\]
这里取对数的是规定单位下的数值, 不是直接对有量纲物理量取对数. 原课程参考通量的 $3.63\times10^{-20}$ 是上述零点的近似.

实际滤镜收一段频率, 并非一个点. 若已把响应及探测方式吸收到归一化权重 $W_X(\nu)$ 中, $\int W_X\dd\nu=1$, 则能量响应约定下定义 $\langle F_\nu\rangle_X=\int F_\nu W_X\dd\nu$, 再代入同一零点比值. 光子计数器的计数率正比于 $\int F_\nu T_X(\nu)\dd\nu/(h\nu)$, 因而应对目标和参考源使用相同的光子权重. 忽略这个区别会把不同光谱误当作相同通量. Vega 与 AB 的通带偏移由参考光谱与响应积分决定, 不是所有滤镜通用的一组常数.

\begin{example}
一个 AB 星等为20的源有 $F_\nu=3631\times10^{-8}\,\mathrm{Jy}=36.31\,\mu\mathrm{Jy}$. 若另一源星等22.5, 其通量为前者的十分之一. 这只比较同一频率或同一响应约定下的通量, 尚未给出两者总光度之比.
\end{example}

绝对星等 $M_X$ 定义为将同一源放到 $10\,\mathrm{pc}$ 时的视星等. 静态无消光条件下, 通量平方反比给出距离模数
\[
\mu\equiv m_X-M_X=5\log_{10}\frac{D}{10\,\mathrm{pc}},
\qquad D=10\,\mathrm{pc}\,10^{\mu/5}.
\]
例如 $\mu=30$ 对应 $D=10\,\mathrm{Mpc}$. 消光使 $m_X$ 增加 $A_X$; 对宇宙学源还需光度距离和把红移光谱映入观测通带的 $K$ 修正. 因而实际使用的是 $m_X=M_X+\mu(D_L)+A_X+K_X$, 并应声明各项是否已校正.

颜色 $m_X-m_Y$ 是两通带通量比的对数. 同一近邻源的距离项抵消, 故它比较光谱形状而不是总光度. 消光一般不抵消, 因为 $A_X\ne A_Y$. 恒星年轻、金属丰度较低或尘埃较少都可能使某些颜色变蓝; 只凭一个颜色不能唯一恢复年龄.

\originalsection{面亮度轮廓与星系形态}
令 $R$ 为投影半径, $I(R)$ 为单位投影面积的光度. 若轴对称, 环带面积为 $2\pi R\dd R$, 故 $L=2\pi\int_0^\infty I(R)R\dd R$. 指数盘有
\[
I(R)=I_0\ee^{-R/R_s},\quad L=2\pi I_0R_s^2,
\quad \frac{L(<R)}L=1-\ee^{-x}(1+x),\quad x=R/R_s.
\]
半光半径 $R_e$ 由右式等于 $1/2$ 确定, $R_e\simeq1.678R_s$. $R_s$ 和 $R_e$ 不能混用.

更一般的 Sérsic 形式可按半光半径写为
\[
I(R)=I_e\exp\left\{-b_n\left[(R/R_e)^{1/n}-1\right]\right\}.
\]
代换 $u=b_n(R/R_e)^{1/n}$ 给出
\[
L=2\pi nI_e\ee^{b_n}R_e^2b_n^{-2n}\Gamma(2n),\qquad
\gamma(2n,b_n)=\tfrac12\Gamma(2n).
\]
第二式才使 $R_e$ 成为半光半径. 指数盘为 $n=1$, $b_1\simeq1.678$; de Vaucouleurs 轮廓为 $n=4$, $b_4\simeq7.669$. 原讲义的中心归一化形式 $I_0\exp[-(R/R_0)^{1/n}]$ 与此式等价, 但 $R_0=R_e/b_n^n$.

\begin{center}
\begin{tikzpicture}
\begin{axis}[width=8.5cm,height=5cm,xlabel={$R/R_e$},ylabel={$I/I_e$},ymode=log,xmin=0.03,xmax=4,ymin=0.002,ymax=300,grid=major,legend style={font=\small}]
\addplot[structurecolor,thick,domain=0.03:4,samples=150]{exp(-1.678*(x-1))};
\addlegendentry{指数盘 $n=1$}
\addplot[black,dashed,domain=0.03:4,samples=150]{exp(-7.669*(x^0.25-1))};
\addlegendentry{$n=4$}
\end{axis}
\end{tikzpicture}
\end{center}
图比较同一 $I_e,R_e$ 下的理想轮廓, 不代表两个模型总光度相同. $n=4$ 的中心更集中、外翼也更宽. 测得一个轮廓并不直接给出质量: 还需质光比及其径向变化.

螺旋星系通常具有恒星和气体盘, 有持续恒星形成, 旋转对支持作用重要; 椭圆星系常以老恒星为主, 随机速度弥散较重要; 透镜状星系有恒星盘而冷气体和恒星形成较弱. 这些是群体上的趋势, 不应把每个椭圆星系都视为完全不转动, 或把形态分类当作唯一的演化时间序列. 环境与并合史也是需要动力学解释的变量.

\originalsection{光度函数及它的积分条件}
按单位体积、单位光度区间定义 $\phi(L)\dd L$. Schechter 经验函数为
\[
\phi(L)\dd L=\phi_*\left(\frac{L}{L_*}\right)^\alpha
\ee^{-L/L_*}\frac{\dd L}{L_*}.
\]
亮端由指数截断, 暗端由 $\alpha$ 控制. 令 $x=L/L_*$, 总数密度与光度密度分别是
\[
n=\phi_*\Gamma(\alpha+1)\quad(\alpha>-1),\qquad
j=\phi_*L_*\Gamma(\alpha+2)\quad(\alpha>-2).
\]
两个收敛条件不同. 课程例子的 $\alpha\simeq-1.09$ 不能直接积分到零光度求有限总数, 却允许有限总光度. 实际样本有选择限, 可写 $n(>L_{\min})=\phi_*\Gamma(\alpha+1,L_{\min}/L_*)$. 此处不把经验暗端外推到零当作无限多真实星系的观测证明.

课程给出的 $\phi_*\sim0.02h^3\,\mathrm{Mpc^{-3}}$, $L_*\sim10^{10}h^{-2}\Ls$ 是相应样本与通带的数量级例子. 这里的 $h$ 由\[
H_0=100h\,\mathrm{km\,s^{-1}Mpc^{-1}}
\]
定义; 改变距离标定同时改变所推断的体积与光度, 因而这两个 $h$ 次幂应一起保留.

\originalsection{速度、引力质量与暗物质线索}
光谱线中心的偏移给出红移, 定义 $1+z=\lambda_{\rm obs}/\lambda_0$. 对低红移近邻样本, 在忽略特有速度影响时, $D\simeq cz/H_0$. 星系内不同位置的谱线还给出视线速度. 倾斜盘的测量近似 $v_{\rm los}=v_c\sin i\cos\varphi$; 倾角不明会影响从投影速度恢复真实旋转速度.

若近似球对称且圆轨道成立, 牛顿引力给出 $v_c^2=GM(<r)/r$. 中心质量外应有 $v_c\propto r^{-1/2}$; 若观测外部旋转曲线近乎平坦, 则
\[
M(<r)\simeq\frac{v_c^2r}{G},\qquad
\rho(r)\simeq\frac{v_c^2}{4\pi Gr^2}.
\]
后一式是球模型下的推断, 不可直接当作薄盘的精确反演. 平坦曲线提示可见光之外有较广的引力质量分布, 但单一观测并不确定粒子种类. 随机运动支持的星系则需速度弥散、轨道各向异性与分布函数, 后文用 Jeans 方程建立这种联系.

原课程按2023年的背景讨论冷、温、自相互作用与极轻玻色子暗物质, 以及修改动力学的不同解释路线. 这些模型对小尺度结构、散射或波动尺度的预测不同. 重子气体、昏暗恒星和已知轻中微子不能在同一组观测限制下无条件承担全部暗物质; 微引力透镜等限制依赖致密天体的质量范围和空间分布. 本章不把课程中的概述改写成任何候选已经被发现或在全部参数区间排除的结论.

\originalsection{恒星种群合成是一个逆问题}
星系光谱由不同恒星的谱叠加, 简化模型写为 $L_\nu=\sum_aN_a\ell_{\nu,a}$. 恒星类别的温度不同, 其连续谱可用黑体作初步比较, 但真实谱含吸收线、金属丰度和表面重力信息. 一个更物理的前向模型为
\[
L_\nu(t)=\int_0^t\psi(t')\,\ell_\nu^{\rm SSP}(t-t',Z(t'))\dd t',
\]
其中 $\psi$ 为恒星形成率, $\ell_\nu^{\rm SSP}$ 是按初始质量函数归一化的单次形成恒星种群光谱, $Z$ 为金属质量分数. 再加入尘埃传输与滤镜响应才得到探测器数据. 不同年龄、金属丰度和尘埃可能产生近似相同颜色, 所以从少数通带恢复完整形成史通常不唯一. 后面的化学演化章会补充独立的守恒条件, 使观测、动力学和演化相互约束.
